\section*{Capítulo \ref{ch:g6:latitude-longitude} --- Latitud, longitud y husos horarios}

\begin{solution}{exo:g6:latitude-longitude:1}
``La \omterm{def:g6:latitude-longitude:latitude}{latitud} se mide desde el \omterm{def:g3:continents-and-oceans:equator}{ecuador} y va de $0^\circ$ a $90^\circ$ (norte
o sur); la \omterm{def:g6:latitude-longitude:longitude}{longitud} se mide desde el \omterm{def:g6:latitude-longitude:longitude}{meridiano de origen} y va de $0^\circ$
a $180^\circ$ (este u oeste)''.
\end{solution}

\begin{solution}{exo:g6:latitude-longitude:2}
El \omterm{def:g3:continents-and-oceans:poles}{Polo Norte}: $90^\circ$\,N. El \omterm{def:g3:continents-and-oceans:poles}{Polo Sur}: $90^\circ$\,S. Una ciudad en el
\omterm{def:g3:continents-and-oceans:equator}{ecuador}: $0^\circ$.
\end{solution}

\begin{solution}{exo:g6:latitude-longitude:3}
Sídney: \omterm{def:g3:continents-and-oceans:equator}{hemisferio sur}, este. Nueva York: \omterm{def:g3:continents-and-oceans:equator}{hemisferio norte}, oeste. El
Cairo: \omterm{def:g3:continents-and-oceans:equator}{hemisferio norte}, este. Río de Janeiro: \omterm{def:g3:continents-and-oceans:equator}{hemisferio sur}, oeste.
\end{solution}

\begin{solution}{exo:g6:latitude-longitude:4}
(a) Verdadero. (b) Falso: los \omterm{def:g6:latitude-longitude:latitude}{paralelos} se acortan hacia los \omterm{def:g3:continents-and-oceans:poles}{polos}; el más
largo es el \omterm{def:g3:continents-and-oceans:equator}{ecuador}. (c) Verdadero. (d) Falso: pasa por Greenwich, en
Londres.
\end{solution}

\begin{solution}{exo:g6:latitude-longitude:5}
El Cairo. Nueva York. Tokio está a unos $36^\circ$\,N, $140^\circ$\,E (un
poco por encima del \omterm{def:g6:latitude-longitude:latitude}{paralelo} $30^\circ$\,N, un poco más allá del \omterm{def:g6:latitude-longitude:longitude}{meridiano}
$120^\circ$\,E, sin llegar a $150^\circ$\,E).
\end{solution}

\begin{solution}{exo:g6:latitude-longitude:6}
$45 \div 15 = 3$: sus mediodías están separados por 3~horas. Es mediodía
primero en la ciudad~B, que está más al este.
\end{solution}

\begin{solution}{exo:g6:latitude-longitude:7}
$73 \div 15 \approx 4.9$, redondeado a 5: el huso es $+5$. Cuando es
mediodía en Greenwich, son las $12 + 5 = 17$, es decir, las 5 de la tarde,
en Bombay.
\end{solution}

\begin{solution}{exo:g6:latitude-longitude:8}
$30 \times 111 = 3\,330$~km. El \omterm{def:g3:continents-and-oceans:poles}{Polo Norte} está a $90^\circ$:
$90 \times 111 = 9\,990$~km, unos 10\,000~km.
\end{solution}

\begin{solution}{exo:g6:latitude-longitude:9}
$9 - (-5) = 14$~horas de diferencia. $8 + 14 = 22$: son las 10 de la noche
del lunes en Tokio.
\end{solution}

\begin{solution}{exo:g6:latitude-longitude:10}
Tiene razón. La Tierra gira de oeste a este, así que las ciudades del este
son llevadas primero hacia el Sol: lo ven salir, llegar a su punto más alto
y ponerse antes que las ciudades situadas más al oeste.
\end{solution}

\begin{solution}{exo:g6:latitude-longitude:11}
(a) $42 \times 3 = 126$~segundos, algo más de 2~minutos. (b)
$126 \div 240 \approx 0.5$: alrededor de medio grado de \omterm{def:g6:latitude-longitude:longitude}{longitud}. (c)
$0.5 \times 111 \approx 55$~km (con $0.525$: unos 58~km). Un error de unos
cincuenta kilómetros: suficiente para no ver una isla, o para encallar en
una \omterm{def:g1:landscapes-around-us:water}{costa} que se creía lejana.
\end{solution}

\begin{solution}{exo:g6:latitude-longitude:12}
(a) $18 - 9 = 9$: las 9 de la mañana del lunes en Greenwich. (b)
$9 + 10 = 19$: las 7 de la tarde del lunes en Greenwich. (c) $19 - 8 = 11$:
las 11 de la mañana del lunes en Los Ángeles. El avión salió a las 6 de la
tarde del lunes (hora de Tokio) y aterriza a las 11 de la mañana del mismo
lunes (hora de Los Ángeles): los relojes de llegada marcan una hora anterior
a la de los relojes de salida, porque el avión cruzó la \omterm{def:g6:latitude-longitude:date-line}{línea de cambio de
fecha} hacia el este y volvió un día atrás.
\end{solution}

\begin{solution}{pb:g6:latitude-longitude:1}
\textbf{1.} Ana (Río de Janeiro, $23^\circ$\,S) y Mia (Auckland,
$37^\circ$\,S).

\textbf{2.} Ana y Sam ($43^\circ$\,O y $118^\circ$\,O). Lena está en el
\omterm{def:g6:latitude-longitude:longitude}{meridiano de origen}.

\textbf{3.} Lena: Londres está a $52^\circ$\,N, más lejos del \omterm{def:g3:continents-and-oceans:equator}{ecuador} que
todas las demás ciudades.

\textbf{4.} Mia: Auckland está a $175^\circ$\,E, a solo $180 - 175 = 5$
grados del \omterm{def:g6:latitude-longitude:longitude}{meridiano} $180^\circ$.

\textbf{5.} $360^\circ \div 24 = 15^\circ$.

\textbf{6.} Londres: $0$. Río de Janeiro: $43 \div 15 \approx 2.9$, al
oeste, así que $-3$. Tokio: $140 \div 15 \approx 9.3$, al este, así que
$+9$. Auckland: $175 \div 15 \approx 11.7$, al este, así que $+12$. Los
Ángeles: $118 \div 15 \approx 7.9$, al oeste, así que $-8$.

\textbf{7.} Río de Janeiro: $20 - 3 = 17$, las 5 de la tarde del viernes.
Tokio: $20 + 9 = 29$; $29 - 24 = 5$: las 5 de la madrugada del sábado.

\textbf{8.} Tokio está muy al este: sus relojes van 9~horas por delante de
los de Londres. A las 8 de la tarde en Londres, en Tokio ya ha pasado la
medianoche, así que su sábado ya ha empezado.

\textbf{9.} $20 - 8 = 12$: las 12 del mediodía del viernes.

\textbf{10.} $12 - (-8) = 20$~horas por delante.

\textbf{11.} Al cruzar la \omterm{def:g6:latitude-longitude:date-line}{línea de cambio de fecha} hacia el este, el barco
vuelve al día anterior.

\textbf{12.} $20 + 12 = 32$; $32 - 24 = 8$: cuando empieza la llamada, son
\textbf{las 8 de la mañana del sábado} en Auckland. Mia se une a la llamada
a la hora del desayuno, mientras Sam come el viernes a mediodía.
\end{solution}
